\section*{Your turn}

Check out tag \code{ch17} of the companion repository and run \code{npm ci}
once. Exercises 1 and 2 compose files and need nothing running. Exercises 3 to 8
need the seven services and the database, so start the stack the way
chapter~\ref{ch:identity-and-authorization} does; the case scripts start their
own routers on ports of their own, so leave the composed stack up rather than
stopping it.

No service changes at this tag, so there is no new Postman collection.
Exercise 8 uses \code{postman/mosaic-router}, which is the one that goes through
the router and has been there since chapter~\ref{ch:enter-the-router}.

\begin{enumerate}
  \item \textbf{Watch the blame move.} From the repository root:

  \begin{minted}[fontsize=\footnotesize,breaklines]{text}
node scripts/satisfiability-cases.mjs --print satisfiability-names-the-first-route
  \end{minted}

  Four compositions, four different root fields named, one fault throughout.
  Before reading the fourth line, predict which field it will name from the
  edit its label describes.

  \item \textbf{Make the composer blame a field you choose.} Open
  \code{schema/catalog.graphql} and swap the declaration order of
  \code{products} and \code{browseProducts} in \code{type Query}. Run exercise 1
  again. The case now fails, and it should: its four expected answers encode the
  committed declaration order, so changing that order is exactly the thing it
  exists to notice. Before running it, predict which of the four runs report a
  different root field than they did in exercise 1. Then put the file back.

  Do not move \code{productById} above \code{browseProducts} for this. The third
  run deletes the block between them by position and will throw rather than
  compose, which tells you nothing about the composer.

  \item \textbf{Read a plan cache key.} Start the stack, then:

  \begin{minted}[fontsize=\footnotesize,breaklines]{text}
node scripts/planner-cases.mjs --print plan-key-ignores-almost-everything
  \end{minted}

  Eleven documents and one number. If the case fails on a document you added
  rather than passing, you have found something this chapter does not describe,
  which is worth more than the exercise.

  \item \textbf{Break the class.} Add a twelfth document to that case's list
  that you expect to share the key, and a pair to the next case that you expect
  to split it. Good candidates for the first: the same query with its fields in
  the same order but on different lines, or with a directive the schema does not
  restrict. Good candidates for the second: a nested field's alias rather than a
  root field's, and a variable with a default value. Two of those four will
  surprise you, and which two is the point.

  \item \textbf{Cost yourself a plan without changing a query.} The case
  scripts start their own routers and remove them again, so to watch this on the
  stack you started, add the debug block to \code{router/config.yaml} and
  restart the router:

  \begin{minted}[fontsize=\footnotesize,breaklines]{yaml}
engine:
  debug:
    enable_cache_response_headers: true
  \end{minted}

  Now send the storefront query twice with its two fields in each order, reading
  the plan cache header each time. Both orders miss on their first send. Then
  count how many orderings a client that builds selection sets from an unordered
  collection could emit for that one query, and hold the answer against
  \code{execution\_plan\_cache\_size}, which defaults to 1024 entries. Decide at
  what number of distinct client queries you would start to care. Take the block
  out again when you are done, or the next chapter's captures will not match.

  \item \textbf{Confirm the observer effect yourself.} Run:

  \begin{minted}[fontsize=\footnotesize,breaklines]{text}
node scripts/planner-cases.mjs --print tracing-never-hits-the-plan-cache
  \end{minted}

  Then read the planning duration on the last line, run it again, and read it
  again. The two numbers differ and neither is the number a production request
  pays. Decide what you would put in an access log instead.

  \item \textbf{Find the duplicate on the wire.} Ask the router for two products
  through \code{nodes}, naming the same identifier twice, with
  \code{X-WG-Trace: true}, and count the representations in the traced input
  against the items in \code{data}. Two go out and three come back. Then post
  the same three representations, duplicate included, straight to Catalog's
  \code{\_entities} on 5101, the way
  chapter~\ref{ch:the-first-cut-extracting-catalog} did, and count what comes
  back. The subgraph answers all three, because the deduplication you just
  watched belongs to the router and Catalog never sees it.

  \item \textbf{Verify it in Postman.} There is no collection for this chapter,
  so use the one that goes through the router:

  \begin{minted}[fontsize=\footnotesize,breaklines]{text}
npx newman run postman/mosaic-router.postman_collection.json \
    -e postman/mosaic-router.local.postman_environment.json
  \end{minted}

  Run it twice without restarting the router. Every request in the second pass
  is a document the router has already planned, and nothing in the collection
  can see that. Now make it see it. Turn on the debug setting that reports cache
  results in response headers:

  \begin{minted}[fontsize=\footnotesize,breaklines]{yaml}
engine:
  debug:
    enable_cache_response_headers: true
  \end{minted}

  Restart the router, then write a test asserting that the plan cache header
  says \code{HIT}. Then decide whether you would ship a collection whose second
  run asserts something its first run cannot, and what that says about running
  it in CI.
\end{enumerate}

If you only have time for one of these, make it exercise 2, because it is the
one that can still surprise somebody who has read the chapter. Predicting which
runs change is easy; predicting the ones that do not is where you find out
whether you believed the mechanism or just remembered the table.

Exercise 5 is the one that will change something you do. A plan cache sized for
your traffic is sized for the number of distinct documents your clients emit,
and that is not the number of queries anybody wrote. The router will not
complain about the difference. It will evict.
